% id: 143-19
% book: 베이직쎈 미적분1 · 09 정적분의 활용
% section: 유형 092 역함수의 그래프와 넓이
% figure: tikz:fig-143-19
% answer: 
% answer_source: 
% variant_level: 0 (원본 전사)
% 이 파일은 build.mjs 가 items.json 에서 만든다. 고칠 때는 items.json 을 고친다.
\smash{\raisebox{1.5mm}{\dmkichul{베이직쎈 미적분1 143쪽 19번}}}
\noindent\begin{minipage}[t]{0.58\linewidth}\vspace{0pt}\raggedright 오른쪽 그림은 함수 $y=f(x)\ (x\ge 1)$와 그 역함수 $y=g(x)$의 그래프이다. 두 곡선 $y=f(x)$, $y=g(x)$가 두 점 $(1{,}\mkern3mu\mathopen{}1)$, $(3{,}\mkern3mu\mathopen{}3)$에서 만나고 $\displaystyle\int_{1}^{3} f(x)\mkern3mu dx=3$일 때, 두 곡선 $y=f(x)$, $y=g(x)$로 둘러싸인 도형의 넓이를 구하시오.\end{minipage}\hfill\begin{minipage}[t]{0.4\linewidth}\vspace{0pt}\centering \resizebox{\ifdim\width>\linewidth\linewidth\else\width\fi}{!}{% 143-19(143쪽) — 함수 $y=f(x)\ (x\ge 1)$ 와 그 역함수 $y=g(x)$ 의 그래프, 직선 $y=x$, 두 곡선이 두 점 $(1,\,1)$, $(3,\,3)$ 에서 만나 둘러싸는 부분에 색칠. 스캔 화질이 낮아 TikZ 로 다시 그림.
% 원문에 있는 것만 옮긴다: 두 곡선 · 직선 $y=x$ · 색칠한 도형 · 안내 점선 · 라벨 O, 1, 3(두 축), x, y, $y=f(x)$, $y=x$, $y=g(x)$. 넓이나 $\int_{1}^{3} f(x)\,dx$ 의 값은 그림에 적지 않는다.
\begin{tikzpicture}[x=0.48cm, y=0.48cm, line width=0.6pt]
  \fill[black!12] plot[domain=1:3, samples=81, smooth] (\x,{1+sqrt(2*((\x)-1))}) -- plot[domain=3:1, samples=81, smooth] (\x,{1+0.5*((\x)-1)*((\x)-1)}) -- cycle;
  \draw[->, >=stealth, line width=0.7pt] (-1.2,0) -- (4.2,0) node[below=1pt, font=\small] {$x$};
  \draw[->, >=stealth, line width=0.7pt] (0,-0.5) -- (0,4.15) node[left=1pt, font=\small] {$y$};
  \draw[densely dashed, line width=0.4pt] (0,3) -- (3,3);
  \draw[densely dashed, line width=0.4pt] (0,1) -- (1,1);
  \draw[densely dashed, line width=0.4pt] (1,0) -- (1,1);
  \draw[densely dashed, line width=0.4pt] (3,0) -- (3,3);
  \draw[line width=0.5pt] (-0.95,-0.95) -- (3.5,3.5);
  \draw[domain=1:3.3, samples=81, smooth] plot (\x,{1+0.5*((\x)-1)*((\x)-1)});
  \draw[domain=1:3.9, samples=81, smooth] plot (\x,{1+sqrt(2*((\x)-1))});
  \node[below right=-1pt and -3pt, font=\small] at (0,0) {$\pt{O}$};
  \node[below=1pt, font=\small] at (1,0) {$1$};
  \node[below=1pt, font=\small] at (3,0) {$3$};
  \node[left=1pt, font=\small] at (0,1) {$1$};
  \node[left=1pt, font=\small] at (0,3) {$3$};
  \node[anchor=east, font=\small] at (3.05,3.85) {$y=f(x)$};
  \node[anchor=west, font=\small] at (3.3,3.9) {$y=x$};
  \node[anchor=west, font=\small] at (3.28,2.95) {$y=g(x)$};
\end{tikzpicture}
}\end{minipage}\par
